\documentclass[11pt,a4paper]{article}
\usepackage{isabelle,isabellesym}

\usepackage[T1]{fontenc}
\usepackage[utf8]{inputenc}

\usepackage{amsmath}
\usepackage{amssymb}

% this should be the last package used
\usepackage{pdfsetup}

\urlstyle{rm}
\isabellestyle{literal}

% Provide any Isabelle symbols not defined by isabellesym.sty.
\providecommand{\isasymto}{\ensuremath{\to}}
\providecommand{\isasymrightarrow}{\ensuremath{\rightarrow}}
\providecommand{\isasymsubseteq}{\ensuremath{\subseteq}}
\providecommand{\isasymmp}{\ensuremath{\mp}}
\providecommand{\isasympm}{\ensuremath{\pm}}
\providecommand{\isasymne}{\ensuremath{\neq}}

% Session / entry names, defined in sessionnames.tex.
% The AFP entry / session names, defined once here and used throughout the
% document, so that a rename is one edit.
\newcommand{\sessMTTM}{\texttt{Multitape\_TM\_Substrate}}
\newcommand{\sessAE}{\texttt{Multitape\_Alphabet\_Enlargement}}
\newcommand{\sessAR}{\texttt{Multitape\_Alphabet\_Reduction}}
\newcommand{\sessART}{\texttt{Multitape\_Alphabet\_Roundtrip}}
\newcommand{\sessNTR}{\texttt{Nondeterministic\_Tape\_Reduction}}
\newcommand{\sessDTR}{\texttt{Deterministic\_Tape\_Reduction}}
\newcommand{\sessNTH}{\texttt{Nondeterministic\_Time\_Hierarchy}}


\begin{document}

% \sloppy: raise the line-breaking tolerance so LaTeX rearranges breaks
% around the long unbreakable \isa{} identifier tokens (in text blocks and
% proof displays) rather than letting them overrun the right page edge.
% Verified to clear every Overfull \hbox >= 20pt in this document.
\sloppy

\title{Multitape Alphabet Enlargement and\\
       the Linear Speedup Theorem\\[6pt]
       \large Formal Proof in Isabelle/HOL}
\author{Andr\'as Z.~Salamon \and Michael Wehar}
\maketitle

\begin{abstract}
Machine-checked formal proof in Isabelle/HOL of the linear
speedup theorem for multitape Turing machines, in both its
deterministic and nondeterministic forms, via the
alphabet-enlargement construction.  The theorem is originally
due to Hartmanis and Stearns~\cite{Hartmanis1965:computational};
the textbook proof followed here is by Hopcroft and
Ullman~\cite[Theorems 12.3 and 12.4]{Hopcroft1979:introduction}
and its nondeterministic corollary.

Given a multitape Turing machine $M$ with alphabet $\Sigma_M$
and a chosen block size $c$, we build an enlarged machine
$M' = \mathrm{alphabet\_enlarge}\, M$.  Its alphabet is
$\Sigma_M^c \cup \{\Box, \rhd\}$ (with a fresh blank and left-end
marker), and it simulates $c$ steps of $M$ at a time, reading its
input in blocks of $c$ symbols.  A separate wrapping step folds the
block-encoder into the machine, so that the textbook statements
compare languages over $M$'s original alphabet.

This session (\sessAE) establishes the following
headline theorems:

\begin{itemize}
\item \emph{Time bound}: if $M$ accepts $w$ in time $T(n)$, then
  $M'$ accepts $\mathrm{encode\_input}\, bl_M\, w$ in time
  $2 \lceil n/c \rceil + 8 \lceil T(n)/c \rceil + 4$ (theorem
  \texttt{alphabet\_enlarge\_time\_explicit}).  Both the input length
  and the running time are divided by the group size $c$.  The same bound
  holds for rejecting runs: if $M$ reaches its reject state on $w$ in time
  $T(n)$, so does $M'$ on the encoded input
  (\texttt{alphabet\_enlarge\_terminal\_time}).
\item \emph{Language biconditional}:
  $\mathrm{encode\_input}\, bl_M\, w \in L(M') \iff w \in L(M)$ for every
  well-formed $M$, with no determinism hypothesis.
\item \emph{Determinism-free building block}: for an arbitrary,
  possibly nondeterministic, well-formed $M$, the language
  biconditional and the time bound hold together
  (\texttt{alphabet\_enlarge\_nae}), with no $\mathrm{det\_ktm}$
  hypothesis.  This is the raw (encoded-input) form over the
  enlarged alphabet; wrapping it in the inline encoder yields the
  same-alphabet speedup corollaries below, each proved in a
  deterministic and a nondeterministic version.
\item \emph{HU 12.4 (linear-$T$ case)}: for $M$ with $k \geq 2$ tapes, as
  in Hopcroft and Ullman, a linear time bound $T(n) \leq d_0 n + b$, and any
  $\varepsilon = 1/q$ with $q$ a positive integer, there exists a
  multitape $M''$ over the original alphabet with $L(M'') = L(M)$ running in
  time $(1 + \varepsilon)|w| + K$ (exactly $|w| + |w| \,\mathrm{div}\, q + K$)
  --- matching Hopcroft and Ullman's $(1 + \varepsilon)n$, with the additive
  constant now an explicit closed expression $K = 28 + 8 d_0 + 8 b$ --- once the
  block size $c$ is taken large enough (it grows with $q$ and the slope
  $d_0$).  It is proved both in a deterministic version ($M$ deterministic
  $\Rightarrow$ $M''$ deterministic) and a nondeterministic version (no
  determinism hypothesis); the latter is the Hopcroft--Ullman corollary
  $\mathrm{NTIME}(cn) = \mathrm{NTIME}((1+\varepsilon)n)$, stated there
  without proof.
\item \emph{HU 12.3 (super-linear-$T$ case)}: for $k \geq 2$ and $T$ satisfying
  the growth hypothesis
  $\forall d. \exists N. \forall n \geq N.\, d \cdot n \leq T(n)$,
  $M''$ runs in time $\varepsilon\, T(|w|) + K$ for
  $|w| \geq N_0$ (the exact integer bound is
  $T(|w|) \,\mathrm{div}\, q + K \leq \varepsilon\, T(|w|) + K$; the
  linear-in-$|w|$ overhead is absorbed into $T(|w|)/(2q)$ via the growth
  hypothesis).  Again this is proved both
  in a deterministic and a nondeterministic version, with the latter giving
  $\mathrm{NTIME}(T(n)) = \mathrm{NTIME}(c \cdot T(n))$ for any
  $c > 0$.
\item \emph{Time-complexity convention}: both linear-$T$ headlines are
  additionally restated in Hopcroft and Ullman's time convention ---
  $\max(n + 2, \cdot)$ here, against their $\max(n + 1, \cdot)$, since the
  substrate reads the left endmarker first --- in three forms, for different
  uses.  The first holds on all inputs and keeps the explicit constant,
  $|w| + |w|\,\mathrm{div}\,q + K$.  The \emph{eventual} form is about the
  same machine and drops the constant, running in exactly
  $|w| + |w|\,\mathrm{div}\,q$ once $|w|$ passes the closed threshold
  $q(q + 1)K$; below it the machine does not meet that bound, and nothing
  forces a valid machine to (the minimal counterexample
  \texttt{clean\_convention\_unattainable}, proved in the substrate entry
  \sessMTTM).  The \emph{table} form adds Hopcroft and Ullman's own device:
  the encoder decides inputs up to the threshold from a finite table on its
  end-of-input transition, within the floor, so the bound
  $|w| + |w|\,\mathrm{div}\,q$ holds on all inputs.  The table exists but is
  not computed, so the first two forms are the effective statements.
\end{itemize}

The language equivalence, the time bound, and the combined building
block above all hold for nondeterministic machines as well as
deterministic ones (they assume no determinism).  Each
Hopcroft--Ullman speedup corollary is proved twice: once for
deterministic machines, where the result is also shown to stay
deterministic, as in the classical statements, and once with no
determinism assumption, giving the corollary that Hopcroft and Ullman
state but do not prove.

The nondeterministic multitape substrate lives in
the sibling \sessMTTM{} session; this session depends
on \sessMTTM{} and \texttt{HOL-Library}.
\end{abstract}

\paragraph*{Acknowledgement.}
This formalisation was developed with proof engineering by
Anthropic's Claude Opus~4.5, 4.6, 4.7, and 4.8 as the work progressed,
and revised with Opus~5.5, together with other Claude family models in
supporting roles.
Claude Code was used to structure the work and for access to Isabelle/HOL.
We also thank Arthur Freitas Ramos, who provided feedback, contributed
several tooling improvements, and made one change in the revision with
OpenAI Codex, driven by OpenAI's GPT-6 and by Claude Opus~4.8.


\tableofcontents

\newpage

\section{Overview}

\subsection{High-level overview}

In one sentence: this session proves the \emph{linear speedup
theorem} for multitape Turing machines, in both its deterministic and
nondeterministic forms.  The theorem says that any machine can be made
to run a constant factor faster.  The idea is to enlarge the machine's
tape alphabet so that a fixed number $c$ of the original cells is
packed into one ``block'', and then to simulate $c$ of the original
steps at a time.

Two constructions do the work.  The first,
$M' = \mathrm{alphabet\_enlarge}\, M$, rebuilds $M$ so that it works
over the block alphabet and simulates $c$ steps of $M$ in each of its
own steps.  Its input must be presented in the same blocked form.  The
second, $M'' = \mathrm{encoding\_wrap}(M', \ldots)$, folds the
block-encoder into the machine, so that the textbook statements can be
made over $M$'s original alphabet, with no encoding visible from
outside.

The headline results are a time bound and a language equivalence for
$M'$ (packaged together, with no determinism assumption, as
$\mathrm{alphabet\_enlarge\_nae}$), together with the two textbook
corollaries over the original alphabet: Hopcroft and Ullman's
Theorem~12.4 (the linear-time case) and Theorem~12.3 (the super-linear
case).  Each corollary is proved twice: once for deterministic
machines (the \texttt{\_dae} version, which also shows the result
stays deterministic, as in the classical statements) and once with no
determinism assumption (the \texttt{\_nae} version).

Three points help a first-time reader.  First, as in the sibling
\sessAR{} session, a machine here may be nondeterministic by default,
since its transition is a relation.  The \texttt{\_nae} theorems are
therefore genuine statements about nondeterministic machines; they are
the $\mathrm{NTIME}$ corollaries
($\mathrm{NTIME}(cn) = \mathrm{NTIME}((1+\varepsilon)n)$ and
$\mathrm{NTIME}(T(n)) = \mathrm{NTIME}(c\,T(n))$) that Hopcroft and
Ullman state without proof.  Second, because the alphabet grows, the
plain machine $M'$ reads its input in blocked form
$\mathrm{encode\_input}\, bl_M\, w$; the wrapped machine $M''$ hides
that encoding, which is why the corollaries can be stated over the
original alphabet.  Third, the block size $c$ can be any fixed number;
the speedup factor is $1 + \varepsilon$ with $\varepsilon = 1/q$, and
each corollary holds once $c$ is large enough relative to $q$.

Every theorem is stated for an arbitrary machine $M$, so each is a
genuine ``for all $M$'' result.  The sibling entry \sessAR{} goes the
other way: it \emph{reduces} the alphabet to a fixed four-symbol set.

% The vocabulary of the substrate shared by the Alphabet Enlargement and
% Alphabet Reduction documents.  Each root.tex inputs it under its own
% Glossary subsection and then adds its transformation's own terms.
\subsection{Glossary}

The machine type is $\mathrm{ktm}$, the ten-component descriptor
$\mathrm{KTM}\ Q\ \Sigma\ \Gamma\ bl\ le\ \delta\ s\ t\ r\ k$: state
set, input alphabet, tape alphabet, blank symbol, left endmarker,
transition relation, start/accept/reject states, and tape count.
Throughout, $M$ is the source machine.

\paragraph{Projections.}  Positional selectors that read one component
out of $M$:
\begin{itemize}
\item \texttt{k\_ktm}\,$M$ --- the tape count $k_M$ (a natural number);
\item \texttt{Sigma\_ktm}\,$M$ --- the input alphabet $\Sigma_M$;
\item \texttt{\isasymGamma\_ktm}\,$M$ --- the tape alphabet $\Gamma_M$;
\item \texttt{bl\_ktm}\,$M$ --- the blank symbol;
\item \texttt{delta\_ktm}\,$M$ --- the transition relation $\delta_M$;
\item \texttt{t\_ktm}\,$M$ --- the accept state.
\end{itemize}

\paragraph{Shared substrate predicates.}
Each entry opens with the gist, then the precise conditions.
\begin{description}
\item[\texttt{valid\_ktm}\,$M$] \emph{``$M$ is a legal machine''.}  The
  structural invariant: $Q$ and $\Gamma$ finite,
  $\Sigma \subseteq \Gamma$, the three distinguished states lie in $Q$,
  blank and endmarker are tape-but-not-input symbols,
  accept $\neq$ reject, $k > 0$, the transition relation is correctly
  typed, the \emph{left-endmarker read discipline} holds (reading the
  endmarker forces re-writing it and never moving off the left edge),
  and every tape beyond index $k$ is frozen blank.
\item[\texttt{well\_formed\_ktm}\,$M$] \emph{a legal machine that also
  starts cleanly and keeps its endmarker.}  An abbreviation:
  \texttt{valid\_ktm}\,$M$ strengthened with \emph{non-degeneracy}
  (start $\neq$ accept, start $\neq$ reject, endmarker $\neq$ blank) and
  the endmarker \emph{write} discipline \texttt{le\_unique}\,$M$.
  Required wherever a run must genuinely start, the endmarker must be
  detectable, and it must never be forged.
\item[\texttt{le\_unique}\,$M$] \emph{the left endmarker is never
  freshly planted.}  A machine property in the same family as
  \texttt{det\_ktm} --- an opt-in structural guarantee a downstream
  construction may rely on.  Precisely: no transition writes the
  endmarker where it was not already reading it
  ($a'\,j = le \Rightarrow a\,j = le$ over $\delta$).  Folded into
  \texttt{well\_formed\_ktm}; the alphabet transformations use it to
  pin the endmarker, and the $(1+\varepsilon)n$ speed-up machine of the
  enlargement entry is the one construction that deliberately forgoes
  it: it writes a fresh endmarker past the input it has consumed, so it
  is \texttt{valid\_ktm} but not \texttt{well\_formed\_ktm}.
\item[\texttt{det\_ktm}\,$M$] \emph{$M$ is deterministic.}  The
  transition relation is a partial function: each (state, read-vector)
  pair has at most one successor triple --- the functionality predicate
  layered on the otherwise relational $\delta$.
\item[$L(M) = \texttt{Lang\_ktm}\,M$] \emph{the language $M$ accepts.}
  All input words $w$ (with $\mathrm{set}\,w \subseteq \Sigma_M$) from
  which some run reaches the accept state; acceptance is
  \emph{existential} over runs --- nondeterministic.
\item[\texttt{accepts\_in\_time\_ktm}\,$M\,w\,t$] \emph{$M$ accepts
  $w$ within $t$ steps.}  Weak time-bounded acceptance: some run of
  length at most $t$ reaches the accept state.
\end{description}

\paragraph{The input lives on a work tape.}  The substrate has
\emph{no} separate read-only input tape: the input word is placed on
tape $0$, which is an ordinary work tape.  This is why
\texttt{valid\_ktm} requires $\Sigma \subseteq \Gamma$ --- every input
symbol occupies a tape cell, so it must also be a tape symbol --- and it
matters to the alphabet transformations, which re-encode the
input \emph{in place} on tape $0$ rather than copying it to a fresh
tape.


\paragraph{Alphabet-enlargement vocabulary.}
\emph{Alphabet enlargement} is our name for the block
construction; in the classical literature it is the alphabet-growing
half of the Hartmanis--Stearns tape-compression / linear-speedup
theorem and carries no standard short name of its own.
\begin{description}
\item[$c = |{'c}|$] the \emph{group size}: the number of original
  cells packed into one block, equal to
  $\mathrm{card}\,(\mathrm{UNIV} :: {'c}\ \mathrm{set})$ for a type
  parameter ${'c} :: \mathrm{enum}$.
\item[$M' = \mathrm{alphabet\_enlarge}\, M$] the enlarged machine,
  over the block alphabet ${'c} \to {'a}$, with state type
  ${'q} \times ({'a},{'c})\ \mathrm{ae\_stage}$; it simulates $c$
  steps of $M$ per $M'$-stage following Hopcroft--Ullman's
  eight-substep scheme.
\item[$\mathrm{ae\_stage}$] the per-state bookkeeping carried by $M'$
  (validation / simulation substep tag, per-tape buffers, head
  destinations).
\item[$\mathrm{encode\_input}\, bl\, w$] the input encoding: it groups
  $w$ into consecutive blocks of $c$ symbols, each packed as a function
  ${'c} \to {'a}$ holding those $c$ symbols, padding the final block
  with the blank $bl$.  The result is
  a list of $\lceil |w| / c \rceil$ blocks.
\item[$M''$, the wrapped machine]
  $\mathrm{encoding\_wrap}\, M'\, (\mathrm{ae\_pack}\, bl_M)\, c\,
  \Sigma_M$: the raw machine $M'$ with an inline block-encoder.  It keeps
  $M$'s $k$ tapes: the encoder copies the input in blocks onto a second
  tape, and the input tape then serves as one of $M'$'s work tapes, with
  a fresh left-end marker $\rhd$ written past the consumed input instead
  of a rewind.  It works over the original-alphabet
  symbols $\mathrm{Raw}\,\Sigma_M$ (plus the encoded working symbols
  $\mathrm{Enc}\,\Gamma_M$); $\mathrm{ae\_pack}$ is the per-block pack
  function matching $\mathrm{encode\_input}$.
\item[$\texttt{Lang\_user\_wrap}\, M''$] the user-facing language of a
  wrapped machine, over the \emph{original} alphabet:
  $\{w \mid \mathrm{map}\ \mathrm{Raw}\ w \in L(M'')\}$.
\item[\texttt{accepts\_in\_time\_user\_wrap}\,$M''\,w\,t$] the
  corresponding user-facing timed acceptance,
  $\texttt{accepts\_in\_time\_ktm}\, M''\, (\mathrm{map}\
  \mathrm{Raw}\ w)\, t$.
\end{description}

Theorem-name suffixes: \texttt{\_nae} marks the determinism-free
(``nondeterministic alphabet enlargement'') version, \texttt{\_dae}
the deterministic specialisation that additionally preserves
determinism.

\paragraph{The slowdown constant.}
The scheme formalised here is Hopcroft and Ullman's eight-substep
simulation~\cite[Theorems 12.3 and 12.4]{Hopcroft1979:introduction} (already
in their 1969 text~\cite[\S7.1]{Hopcroft1969:formal}), whose per-group cost
is the factor $8$ in the time bound above.  The precise constant is a matter
of presentation, not of the theorem: Balc\'azar, D\'iaz, and
Gabarr\'o~\cite{Balcazar1988:structural} and
Papadimitriou~\cite{Papadimitriou1994:computational} give a factor-$6$
simulation, while van Emde Boas's survey~\cite{vanEmdeBoas1990:machine} reads
the original Hartmanis--Stearns theorem~\cite{Hartmanis1965:computational} as
intending an amortised, non-oblivious simulation that reaches
$(1 + \varepsilon)$ directly rather than through any fixed multiplicative
constant.  Whichever the constant, the group size $c$ divides it away, so we
formalise the scheme cleanest to verify rather than the one with the
smallest factor.  For deterministic machines the floor of $n$ is not
reached:
Rosenberg~\cite[Proposition~2]{Rosenberg1967:realtime} gives a language
accepted in time $2n$, hence in time $(1 + \varepsilon)n$ for every
$\varepsilon > 0$, that no deterministic real-time multitape machine
accepts, so the $(1 + \varepsilon)n$ of Theorem~12.4 cannot in general be
improved to $n$.  For nondeterministic machines it can, by a different
argument: Book and Greibach~\cite{Book1970:quasi-realtime} show that
nondeterministic linear time and real time have the same power.
These are results about on-line machines with a separate input tape; the
construction here makes no claim at the floor itself.

\paragraph{Related AFP developments.}
To our knowledge the Hartmanis--Stearns linear speedup theorem has not been
formalised in any proof assistant, so this entry has no direct predecessor;
the neighbouring Isabelle/HOL developments prove different theorems.
Balbach's \texttt{Cook\_Levin}~\cite{Balbach2023:Cook} shares the model ---
multitape machines with the tape count a value --- and, like this entry,
states a machine's language and its running-time bound as separate theorems,
but treats deterministic polynomial-time reductions for
$\mathrm{NP}$-completeness, not a speedup.  Xu, Zhang, Urban, and Joosten's
\texttt{Universal\_Turing\_Machine}~\cite{Xu2019:universal} is a single-tape
universal machine over a fixed binary alphabet for computability, with no
timing analysis.  The substrate itself descends from Dalvit and Thiemann's
\texttt{Multitape\_To\_Singletape\_TM}~\cite{Dalvit2022:verified}, but makes
the number of tapes a machine's \emph{data} rather than a type parameter, so a
single speedup theorem here ranges over machines of every tape count $k \geq 2$.

\subsection{Scope and limitations}

Three caveats orient a reader deciding whether these results apply.
\begin{itemize}
\item \emph{Original-alphabet statements need the wrap.}  The raw
  enlarged machine $M'$ works on block-encoded input
  ($\mathrm{encode\_input}\, bl_M\, w$); the same-alphabet corollaries
  hold for the \emph{wrapped} machine $M''$, which folds the encoder
  in.  Statements about $M'$ alone carry the visible encoding.
\item \emph{The corollaries need a large enough group.}  The speedup
  factor is $1 + \varepsilon$ with $\varepsilon = 1/q$, and each
  corollary holds only once the group size $c$ is large enough
  relative to $q$ ($q\,(3 + 8 d_0) \leq c$ for the linear
  case, $16 q \leq c$ for the super-linear case).
\item \emph{The super-linear case needs a growth hypothesis.}
  HU~12.3 assumes
  $\forall d.\ \exists N.\ \forall n \geq N.\ d\,n \leq T(n)$ and holds
  for inputs $|w| \geq N_0$; the linear case HU~12.4 instead assumes
  $T(n) \leq d_0 n + b$.
\end{itemize}

\subsection{The Hopcroft--Ullman statements}

So a reader can check this session against the textbook results it
claims, we reproduce the source statements verbatim from Hopcroft and
Ullman~\cite[\S12.2, pp.~290--291]{Hopcroft1979:introduction}, with the
two conventions needed to read them.

\paragraph{Two reading conventions.}
First, in these quoted statements $c$ is Hopcroft and Ullman's
\emph{time-complexity coefficient}; it is \emph{unrelated} to our group
size $c$ (the enlargement block size of the vocabulary above).  Second,
they use the convention that $cT(n)$ abbreviates
$\max(n + 1, \lceil cT(n) \rceil)$: the $n + 1$ floor is the cost of
reading the input and reaching the blank past its end, and it is what
lets the bound hold for \emph{all} $n$ --- the finitely many short
inputs are recognised by finite control in $n + 1$ moves, and only the
long inputs run the simulation.

\paragraph{Theorem 12.3 (linear speed-up).}
\begin{quote}
If $L$ is accepted by a $k$-tape $T(n)$ time-bounded Turing machine
$M_1$, then $L$ is accepted by a $k$-tape $cT(n)$ time-bounded TM $M_2$
for any $c > 0$, provided that $k > 1$ and
$\inf_{n \to \infty} T(n)/n = \infty$.
\end{quote}

\paragraph{Theorem 12.4 (linear-$T$ case).}
\begin{quote}
If $L$ is accepted by a $k$-tape $cn$ time-bounded TM, for $k > 1$ and
for some constant $c$, then for every $\varepsilon > 0$, $L$ is accepted
by a $k$-tape $(1 + \varepsilon)n$ time-bounded TM.
\end{quote}

\paragraph{Corollaries (p.~291).}
\begin{quote}
\emph{To Theorem 12.3:} if $\inf_{n \to \infty} T(n)/n = \infty$ and
$c > 0$, then $\mathrm{DTIME}(T(n)) = \mathrm{DTIME}(cT(n))$.

\smallskip\noindent
\emph{To Theorem 12.4:} if $T(n) = cn$ for some $c > 1$, then
$\mathrm{DTIME}(T(n)) = \mathrm{DTIME}((1 + \varepsilon)n)$ for any
$\varepsilon > 0$.

\smallskip\noindent
\emph{Of Theorems 12.3 and 12.4:}
\textbf{(a)} if $\inf_{n \to \infty} T(n)/n = \infty$, then
$\mathrm{NTIME}(T(n)) = \mathrm{NTIME}(cT(n))$ for any $c > 0$;
\textbf{(b)} if $T(n) = cn$ for some constant $c$, then
$\mathrm{NTIME}(T(n)) = \mathrm{NTIME}((1 + \varepsilon)n)$ for any
$\varepsilon > 0$.
\end{quote}

\paragraph{Correspondence.}
The session's four headline theorems formalise these results over the
substrate:
\begin{description}
\item[\texttt{linear\_speedup\_HU\_12\_3\_dae}] Theorem 12.3, the
  deterministic linear speed-up (super-linear $T$).
\item[\texttt{linear\_speedup\_HU\_12\_3\_nae}] Corollary (a), its
  nondeterministic form.
\item[\texttt{linear\_speedup\_HU\_12\_4\_dae}] Theorem 12.4, the
  deterministic linear-$T$ speed-up.
\item[\texttt{linear\_speedup\_HU\_12\_4\_nae}] Corollary (b), its
  nondeterministic form.
\end{description}
The \texttt{DTIME}/\texttt{NTIME} corollaries are the complexity-class
readings of the deterministic and nondeterministic forms respectively.
The precise formal statements of all four --- read against the
$\max(n + 1, \lceil \cdot \rceil)$ convention above --- follow next.

\subsection{Details of theorem statements}

We first state the three properties of the raw enlarged machine $M'$,
then their determinism-free bundle, then the two textbook corollaries
over the original alphabet.  The only hypothesis on $M$ for the raw
results is \texttt{well\_formed\_ktm}: the time and language
results carry \emph{no} determinism hypothesis, and the deterministic
corollaries add \texttt{det\_ktm}.

\paragraph{\texttt{alphabet\_enlarge\_wf}.}
\[
  \mathrm{valid\_ktm}\, M \ \Longrightarrow\ \mathrm{valid\_ktm}\, M'.
\]
The enlarged machine is again a legal machine.  Unlike reduction, no
alphabet-size side condition is needed: enlargement always yields a
well-typed block machine.

\paragraph{\texttt{alphabet\_enlarge\_language}.}
\[
  \begin{aligned}
    &\mathrm{well\_formed\_ktm}\, M \ \Longrightarrow\
      \forall w.\ \mathrm{set}\, w \subseteq \Sigma_M \longrightarrow {} \\
    &\quad \bigl(\mathrm{encode\_input}\, bl_M\, w \in L(M')\bigr)
      = \bigl(w \in L(M)\bigr).
  \end{aligned}
\]
The language biconditional for the raw machine, with no determinism
hypothesis: $M'$ accepts the block-encoding of $w$ exactly when $M$
accepts $w$.

\paragraph{\texttt{alphabet\_enlarge\_time\_explicit} and
  \texttt{alphabet\_enlarge\_time}.}
Under $\mathrm{well\_formed\_ktm}\, M$, for every input word $w$:
\[
  \begin{aligned}
    &\mathrm{accepts\_in\_time\_ktm}\, M\, w\, (T(|w|)) \longrightarrow {} \\
    &\quad \mathrm{accepts\_in\_time\_ktm}\, M'\,
      (\mathrm{encode\_input}\, bl_M\, w)\, \\
    &\qquad \bigl(2\,\lceil |w|/c\rceil + 8\,\lceil T(|w|)/c\rceil + 4\bigr).
  \end{aligned}
\]
The $\lceil\,\cdot\,/c\rceil$ factors are the source of the speedup:
both the input length and the running time are divided by the group
size $c$.  The coefficients $2$, $8$, $4$ are absolute constants (the
$8$ counts the eight substeps simulated per group); none of them
depends on $M$, on $w$, or even on $c$.  The explicit-constant
statement is \texttt{alphabet\_enlarge\_time\_explicit}; the
\texttt{obtains}-corollary \texttt{alphabet\_enlarge\_time} repackages
it existentially as
$\exists \alpha, f.\ \alpha\,\lceil |w|/c\rceil + 8\,\lceil T(|w|)/c\rceil + f$,
instantiated at $\alpha = 2$, $f = 4$.

Both are the accepting instance of
\texttt{alphabet\_enlarge\_terminal\_time}, which bounds the time to reach
either terminal state ($\mathrm{reaches\_state\_in\_time\_ktm}$): a run of
$M$ that rejects within $T(|w|)$ steps is simulated by a rejecting run of
$M'$ within the same bound.  The encoding wrap keeps this, since its reject
state is the simulated machine's (\texttt{wrap\_reaches\_state}), so a
bound on the time to reject, which a class of decided languages needs, can
be had from the same simulation.

\paragraph{\texttt{alphabet\_enlarge\_nae}.}
The determinism-free building block: under
$\mathrm{well\_formed\_ktm}\, M$ alone, the language biconditional
\emph{and} the same time bound hold together, packaged existentially
($\exists \alpha, f$, with the witnesses $\alpha = 2$, $f = 4$).  This
is the raw, encoded-input statement over the enlarged alphabet; wrapping it in the inline encoder yields the same-alphabet
corollaries below.

\paragraph{\texttt{linear\_speedup\_HU\_12\_4\_nae} /
\texttt{\_dae} (linear-$T$ case).}
Assume $\mathrm{well\_formed\_ktm}\, M$, $k \geq 2$, a linear time bound
$T(n) \leq d_0 n + b$, and $q > 0$, $c > 0$.  Then $M''$ is valid,
$\texttt{Lang\_user\_wrap}\, M'' = L(M)$
(same language, \emph{over the original alphabet}), and, provided
$c$ is large enough ($q\,(3 + 8 d_0) \leq c$), for every
$w$,
\[
  \begin{aligned}
    &\mathrm{accepts\_in\_time\_ktm}\, M\, w\, (T(|w|)) \longrightarrow {} \\
    &\quad \texttt{accepts\_in\_time\_user\_wrap}\, M''\, w\, \\
    &\qquad \bigl(|w| + |w| \,\mathrm{div}\, q + K\bigr), \\
    &\qquad\quad K = 28 + 8 d_0 + 8 b .
  \end{aligned}
\]
The additive constant $K$ is a \emph{literal} closed expression in the
input machine's linear-bound constants $d_0, b$.  The running time is
$|w|\,(1 + 1/q) + K = (1+\varepsilon)\,|w| + K$
with $\varepsilon = 1/q$, the classical linear speedup.  The
\texttt{\_dae} version adds the hypothesis $\mathrm{det\_ktm}\, M$ and
the extra conclusion $\mathrm{det\_ktm}\, M''$ (determinism
preserved); the \texttt{\_nae} version drops both and is the
Hopcroft--Ullman corollary
$\mathrm{NTIME}(cn) = \mathrm{NTIME}((1+\varepsilon)n)$.

\paragraph{\texttt{linear\_speedup\_HU\_12\_3\_nae} /
\texttt{\_dae} (super-linear-$T$ case).}
Assume $\mathrm{well\_formed\_ktm}\, M$, $k \geq 2$, the growth hypothesis
$\forall d.\ \exists N.\ \forall n \geq N.\ d\,n \leq T(n)$, $q > 0$,
and $16 q \leq c$.  Then there are constants $K, N_0$ with $M''$
valid, $\texttt{Lang\_user\_wrap}\, M'' = L(M)$, and for every $w$ with
$|w| \geq N_0$,
\[
  \begin{aligned}
    &\mathrm{accepts\_in\_time\_ktm}\, M\, w\, (T(|w|)) \longrightarrow {} \\
    &\quad \texttt{accepts\_in\_time\_user\_wrap}\, M''\, w\,
      \bigl(T(|w|) \,\mathrm{div}\, q + K\bigr).
  \end{aligned}
\]
For super-linear $T$ the growth hypothesis absorbs the linear-in-$|w|$
overhead into $T(|w|)/(2q)$, leaving running time $T(|w|)/q + K$, a
speedup by an arbitrary factor $q$.  As before, the \texttt{\_dae}
version adds $\mathrm{det\_ktm}$ preservation, and the \texttt{\_nae}
version is the corollary
$\mathrm{NTIME}(T(n)) = \mathrm{NTIME}(c\,T(n))$ for any $c > 0$.

\paragraph{Convention restatements
(\texttt{linear\_speedup\_HU\_12\_4\_nae\_conv} /
\texttt{\_eventual} / \texttt{\_conv\_table}, and the \texttt{\_dae}
versions).}
The four headlines above bound each accepting run with an additive constant
$K$ that depends on $M$'s time bound but not on the input.  Read in Hopcroft and Ullman's $\max(n + 1, \cdot)$ time
convention --- here $\max(n + 2, \cdot)$, since the substrate reads the
left endmarker before the first input symbol --- the linear-$T$ case has
three forms, which suit different uses.  The \texttt{\_conv} form is the
raw bound restated: the wrap accepts every $w \in L(M)$ within
$\max(|w| + 2,\ |w| + |w|\,\mathrm{div}\,q + K)$, provided
$q(3 + 8 d_0) \leq c$, with every number explicit; it is the form to
compose with other constructions.  The \texttt{\_eventual} form is
about the wrap itself: it runs in exactly $|w| + |w|\,\mathrm{div}\,q$
once $|w|$ reaches the closed threshold $q(q + 1)K$, with the same literal
$K = 28 + 8 d_0 + 8 b$ (proved by running the raw bound one denominator
tighter, at $q + 1$, then absorbing $K$ into the extra
$\mathrm{div}$-slack).  Below the threshold the wrap alone cannot meet the
bound: its setup phases overshoot the floor on tiny inputs, and nothing
forces a valid machine to meet it --- the substrate entry \sessMTTM{} proves
\texttt{clean\_convention\_unattainable}, a valid deterministic machine that
accepts its language yet needs three steps on the empty input, one past the
floor $2$.  The counterexample is about one machine, not a lower bound on
its language.

The \texttt{\_conv\_table} form drops the constant on \emph{all} inputs:
$\mathrm{time\_bounded\_conv}\ P\ (\lambda n.\ n + n\,\mathrm{div}\,q)$, so
$P$ accepts every $w \in L(M)$ within
$\max(|w| + 2,\ |w| + |w|\,\mathrm{div}\,q)$.  The machine $P$ is the wrap
with a \emph{first-block patch}
(\texttt{wrap\_prefix\_patch}, theory \texttt{Wrap\_Prefix\_Patch}): while
the encoder packs the first block of $c$ input symbols, a flag records that
no block has yet been closed, and if the input ends inside that block the
encoder's end-of-input transition goes straight to the verdict of a table
of the words of length at most the threshold.  This is Hopcroft and
Ullman's own device: their proof recognises short words in the finite
control while the input is read~\cite[p.~291]{Hopcroft1979:introduction}.
A short input is then decided in $|w| + 2$ steps, exactly the floor, and a
long one runs as the wrap does.  An outer cleanup that rescans the input
(\texttt{finite\_patch} in \sessMTTM) would not do: the rescan costs time
proportional to the cutoff, which here grows with $q$.  The width
hypothesis $q(q+1)K < c$ keeps the threshold inside the first block.  The
table is $L(M)$ restricted to short words, a finite set, so the
\texttt{\_conv\_table} forms assert that the machine exists rather than
construct it; this is the form for class equalities, where existence is
all that is needed, and the \texttt{\_conv} and \texttt{\_eventual} forms
are the effective statements.  The patch preserves determinism, so
\texttt{\_dae\_conv\_table} keeps $\mathrm{det\_ktm}$.
The patch's correctness needs no determinism: the wrap is deterministic
until it starts simulating $M$, which is enough to force its first block.

The superlinear case has no threshold to handle inside the first block:
the growth of $T$ dominates the cost of the outer cleanup, so the
convention forms \texttt{linear\_speedup\_HU\_12\_3\_nae\_conv} and
\texttt{\_dae\_conv} are clean over \texttt{finite\_patch}.

\paragraph{Speedup above a linear band.}
Theory \texttt{Wrap\_Band\_Speedup} states the first-block construction for
an arbitrary time bound $t$ that lies, from some $N_0$ on, above the line
$n + n\,\mathrm{div}\,q$ for a positive $q$.  If a well-formed machine $M$
with at least two tapes accepts each word of its language within
$a\,t(n) + e$ steps, for constants $a$ and $e$, then the patched wrap accepts
the same language within the convention bound $\max(n + 2,\ t(n))$
(\texttt{wrap\_prefix\_speedup\_band}).  The constant factor $a$ is absorbed
by taking the group size large enough, as stated in that theorem; the short
words decided from the table are those of length at most
$\max(N_0,\ q(81 + 128 e))$, and the group size must exceed that length.
The accompanying \texttt{wrap\_prefix\_speedup\_contract} gives validity, the
language (each word mapped through $\mathrm{Raw}$) and, for deterministic
$M$, determinism.

Like the headlines, this theorem bounds accepting runs, which is what a
language class needs.  Rejecting runs are carried over on long words, where
the patched machine runs as the wrap does (\texttt{wrap\_reaches\_state}
above), but not on short ones: there it stops with the table's verdict, so
for a nondeterministic $M$ that both accepts and rejects a short word along
different runs, the rejecting runs are not reproduced.  A bound on the time
to reject therefore combines the two cases for the machine at hand.

\newpage

% Isabelle's auto-generated session.tex contains the correct
% \input{} sequence in declaration order from this session's ROOT
% `theories` block.  Self-scoped: documents only Multitape_Alphabet_Enlargement's
% own theories (the parent Multitape_TM_Substrate is documented separately).
\input{session}

\bibliographystyle{abbrv}
\bibliography{root}

\end{document}
